\section{The Euclidean Algorithm}

\begin{theorem}[Euclidean Algorithm]
Let $a, b \in \mathbb{N} : a > 0, b > 0 \text{ and } a \geq b$. The following algorithm computes $\gcd(a, b)$ in a finite number of steps:
\begin{enumerate}
    \item Let $r_0 := a \text{ and } r_1 := b$.
    \item Set $i = 1$.
    \item Divide $r_{i-1}$ by $r_i$ to obtain a quotient $q_i$ and remainder $r_{i+1}$: $r_{i-1} = r_i q_i + r_{i+1}$.
    \item If $r_{i+1} = 0$, then $r_i = \gcd(a, b)$. The algorithm terminates.
    \item Otherwise, increment $i \rightarrow i+1$ and return to Step 3.
\end{enumerate}
\end{theorem}
\begin{proof}
Let the sequence of remainders be $r_0, r_1, \dots, r_n, r_{n+1}$ where $d = r_n$ is the last non-zero remainder and $r_{n+1} = 0$. We must demonstrate two facts: that $d$ is a common divisor of $a$ and $b$, and that it is the greatest common divisor.

\subsubsection*{a. $d \mid a$ and $d \mid b$}
We proceed by working backward from the end of the algorithm.
\begin{itemize}
    \item The final division step is $r_{n-1} = r_n q_n + r_{n+1}$. Since $r_{n+1} = 0$ and $d = r_n$, this equation becomes $r_{n-1} = d q_n$. This demonstrates $d \mid r_{n-1}$.
    \item The preceding step is $r_{n-2} = r_{n-1} q_{n-1} + r_n$. Since $d \mid r_n$ and we have shown $d \mid r_{n-1}$, it follows that $d$ must divide their linear combination, $d \mid r_{n-2}$.
\end{itemize}
Continuing this backward induction process, we establish that $d$ divides every preceding remainder. This concludes that $d \mid r_1$ (where $r_1 = b$) and $d \mid r_0$ (where $r_0 = a$). Thus, $d$ is a common divisor.

\subsubsection*{b. $d$ is the Greatest Common Divisor}
Let $c \in \mathbb{Z}$ be any common divisor of $a$ and $b$. We work forward down the chain of remainders.
\begin{itemize}
    \item Since $c \mid a$ ($r_0$) and $c \mid b$ ($r_1$), and $r_2 = r_0 - r_1q_1$, it must be that $c \mid r_2$.
    \item Since $c \mid r_1$ and $c \mid r_2$, and $r_3 = r_1 - r_2q_2$, it must be that $c \mid r_3$.
\end{itemize}
Continuing this process, we eventually find that $c$ must divide the last non-zero remainder, $r_n$, so $c \mid d$.
Since $c \mid d$, it is a consequence of divisibility that $c \leq d$. This holds $\forall c$ common divisor, proving that $d$ is the greatest common divisor. \qedhere
\end{proof}

\begin{center}
    \vspace{10pt}
    \line(1,0){450}
    \vspace{25pt}
\end{center}

\begin{proposition}[Extended Euclidean Algorithm]
Let $a, b \in \mathbb{N}$ such that $a \neq 0$ or $b \neq 0$. If $d = \gcd(a, b)$, then $\exists u, v \in \mathbb{Z} : d = au + bv$.
\end{proposition}

\begin{proof}
    The proof relies on the standard Euclidean Algorithm applied to $a$ and $b$, which generates a sequence of remainders $r_i$:
    \begin{align*}
        a &= q_1 b + r_1, \quad \text{with } 0 \le r_1 < b \\
        b &= q_2 r_1 + r_2, \quad \text{with } 0 \le r_2 < r_1 \\
        r_1 &= q_3 r_2 + r_3, \quad \text{with } 0 \le r_3 < r_2 \\
        &\vdots \\
        r_{n-3} &= q_{n-1} r_{n-2} + r_{n-1}, \quad \text{with } 0 \le r_{n-1} < r_{n-2} \\
        r_{n-2} &= q_n r_{n-1} + 0.
    \end{align*}
    
    Since $r_{n-1}$ is the last non-zero remainder, we have $d = \gcd(a, b) = r_{n-1}$.
    
    We proceed by **backward substitution** (as indicated by the notes on the right):
    \begin{itemize}
        \item From the second-to-last equation, $r_{n-1}$ can be expressed as a linear combination of $r_{n-3}$ and $r_{n-2}$:
        $$d = r_{n-1} = r_{n-3} - q_{n-1} r_{n-2}$$
        \item Since $r_{n-2}$ is a linear combination of $r_{n-4}$ and $r_{n-3}$, substituting this expression for $r_{n-2}$ into the equation for $d$ will express $d$ as a linear combination of $r_{n-4}$ and $r_{n-3}$.
        \item By continuing this backward substitution process (i.e., replacing each remainder $r_i$ with its expression from the previous division step in the Euclidean algorithm), we progressively express $d$ as a linear combination of the previous two terms in the remainder sequence:
        $$r_{i} = \text{lin. comb. of } r_{i-2} \text{ and } r_{i-1}$$
        \item This process continues until we reach the first two equations, where $d$ will be expressed as a linear combination of $a$ and $b$:
        $$d = r_{n-1} = au + bv$$
        for some integers $u, v \in \mathbb{Z}$.
    \end{itemize}
\end{proof}

\begin{center}
    \vspace{10pt}
    \line(1,0){450}
    \vspace{25pt}
\end{center}

\newpage
\begin{proposition}[Solvability of Linear Diophantine Equations]
    Let $a, b \in \mathbb{Z}$ be integers, not both zero. The equation $au + bv = c$ has integer solutions $u, v \in \mathbb{Z}$ if and only if $\gcd(a, b) \mid c$.
\end{proposition}

\begin{proof}
    \textbf{$(\impliedby)$ If $\gcd(a, b) \mid c$, solutions exist} \\
    Assume $c = k \cdot \gcd(a, b)$ for some $k \in \mathbb{Z}$. By Bézout's Identity, $\exists u_0, v_0 \in \mathbb{Z}$ so that $au_0 + bv_0 = \gcd(a, b)$.
    
    Multiplying by $k$ yields:
    $$k(au_0 + bv_0) = k \cdot \gcd(a, b) \implies a(ku_0) + b(kv_0) = c$$
    
    Thus, $u = ku_0$ and $v = kv_0$ are integer solutions.

    \textbf{$(\implies)$ If solutions exist, then $\gcd(a, b) \mid c$} \\
    Assume $au + bv = c$ has integer solutions $u, v \in \mathbb{Z}$.
    
    Let $d = \gcd(a, b)$. Since $d \mid a$ and $d \mid b$, $d$ must divide any linear combination of $a$ and $b$:
    $$d \mid (au + bv) \implies d \mid c$$
\end{proof}

\begin{center}
    \vspace{10pt}
    \line(1,0){450}
    \vspace{25pt}
\end{center}
\begin{proposition}[All Integer Solutions to $au+bv=c$]
    Let $a, b \in \mathbb{Z}$ be relatively prime (i.e., $\gcd(a, b) = 1$), and let $u_0, v_0 \in \mathbb{Z}$ be a particular solution so that $au_0 + bv_0 = c$. Then, all integer solutions $(u, v)$ to $au + bv = c$ are of the form:
    $$u = u_0 + kb \quad \text{and} \quad v = v_0 - ka$$
    for any $k \in \mathbb{Z}$.
\end{proposition}

\begin{proof}
    \textbf{Part 1: Verification (The form gives solutions)} \\
    Substituting the general form into the equation:
    \begin{align*}
        a(u_0 + kb) + b(v_0 - ka) &= au_0 + akb + bv_0 - bka \\
        &= (au_0 + bv_0) + (akb - bka) \\
        &= c + 0 \\
        &= c
    \end{align*}
    Thus, $\forall k \in \mathbb{Z}$, the formula yields a solution.

    \textbf{Part 2: Completeness (All solutions must be of this form)} \\
    Assume $(u, v)$ is any arbitrary solution, so $au + bv = c$. We also have the particular solution $au_0 + bv_0 = c$.
    
    Subtracting the two equations gives:
    $$a(u - u_0) + b(v - v_0) = 0 \implies a(u - u_0) = -b(v - v_0) \implies a(u - u_0) = b(v_0 - v)$$
    
    Since $a \mid b(v_0 - v)$ and we are given $\gcd(a, b) = 1$, by Euclid's Lemma, it must be that $a$ divides the other factor, $(v_0 - v)$:
    $$a \mid (v_0 - v)$$
    
    Let $k \in \mathbb{Z}$ be the integer so that $v_0 - v = ka$. Rearranging this gives the form for $v$:
    $$v = v_0 - ka$$
    
    Substituting $v_0 - v = ka$ back into the equality $a(u - u_0) = b(v_0 - v)$:
    $$a(u - u_0) = b(ka) = a(kb)$$
    
    Since $a \ne 0$, we can divide by $a$ to get:
    $$u - u_0 = kb$$
    
    Rearranging this gives the form for $u$:
    $$u = u_0 + kb$$
    
    Hence, all solutions are of the specified form.
\end{proof}

\begin{center}
    \line(1,0){450}
\end{center}